\documentclass[10pt]{article}
\usepackage{braket}
\usepackage{soul}
\usepackage{float}
\usepackage{tabularx}
\usepackage{mathtools}
\usepackage{amsmath, amssymb, amsthm}
\usepackage{tikz}
\usetikzlibrary{arrows}
\usetikzlibrary{decorations.text}
\usepackage{enumitem}
\usepackage{geometry}
\usepackage{fancyhdr}
\usepackage{titlesec}
\usepackage{xltabular}
\usepackage{color}
\usepackage{hyperref}
\geometry{margin=1in}

\pagestyle{fancy}
\fancyhf{}
\cfoot{\thepage}
\title{MATH 2R03 - Assignment 2}
\author{Matthew Farah, \href{mailto:farahm11@mcmaster.ca}{$\langle$farahm11@mcmaster.ca$\rangle$}, 400468251}
\date{\today}

\hypersetup{
	colorlinks=true,
	linkcolor=blue,
	filecolor=magenta,
	urlcolor=blue,
	citecolor=blue,
	anchorcolor=blue
}

\fancyhead{}
\fancyhead[L]{Matthew Farah, \href{mailto:farahm11@mcmaster.ca}{$\langle$farahm11@mcmaster.ca$\rangle$}, 400468251}
\fancyhead[R]{MATH 2R03 - Assignment 2}

\DeclareMathOperator{\spn}{span}
\DeclareMathOperator{\degree}{deg}

\AtBeginEnvironment{align}{\setcounter{equation}{0}}

\newcommand{\comment}[1]{\parbox[t]{0.45\textwidth}{\raggedright #1}}

\setlength{\parindent}{0cm}

\setcounter{tocdepth}{3}
\hypersetup{bookmarksopen=true, bookmarksopenlevel=2, bookmarksdepth=3}

\newcounter{chapter}
\renewcommand{\thechapter}{\Roman{chapter}}

\newenvironment{chaptertitle}{
	\newpage
	\refstepcounter{chapter}
	\begin{center}
	\vspace*{.5em}
	\par
	\Large
	Part \thechapter
	\phantomsection
	\addcontentsline{toc}{section}{\protect{\large{Part \thechapter}}}
}{
	\end{center}
	\vspace{2.5em}
}

\newcounter{question}
\renewcommand{\thequestion}{\arabic{question}}

\newenvironment{question}{
	\refstepcounter{question}
	\phantomsection
	\addcontentsline{toc}{subsection}{\protect{\large{Question \thequestion}}}
	\vspace{1em}
	\par
	\large\textbf{Question \thequestion}
	\vspace{.2cm}
	\par
	\setcounter{subquestion}{0}
	\setcounter{step}{0}
}{\vspace{.5em}}

\newenvironment{solution}{
	\vspace{1em}\par\large\textbf{{Solution:}}\par\vspace{1em}
}{
	\vspace{1em}
	\newpage
}

\newcounter{subquestion}
\renewcommand{\thesubquestion}{\alph{subquestion}}

\newenvironment{subquestion}{
	\refstepcounter{subquestion}
	\phantomsection
	\vspace{1em}\par\large\textbf{(\thesubquestion)}\hspace{-.5em}
	\setcounter{step}{0}
	\addcontentsline{toc}{subsubsection}{\protect{\large{Part \thequestion.\thesubquestion}}}
}{
	\par
	\vspace{1em}
}

\makeatother

\makeatletter

\newcounter{step}
\renewcommand{\thestep}{\Roman{step}}

\newenvironment{step}{
	\refstepcounter{step}
	\vspace{1.5em}\par\noindent\large\textbf{(\thestep)}
}{
	\par
	\vspace{1.5em}
}

\makeatother

\newcolumntype{Y}{>{\centering\arraybackslash}X}
\renewcommand{\arraystretch}{1.8}

\fontsize{10pt}{14pt}\selectfont

\begin{document}

\maketitle

\tableofcontents
\newpage

\begin{question}
	Let $V$ be a vector space over a field $\mathbb{F}$, and let $v_1, \ldots, v_m$ be a list of vectors in $V$. For $k = 1, \ldots, m$ set $w_k \coloneq v_1 + \ldots + v_k$. Show that $\spn(w_1, \ldots, w_{m}) = \spn(v_1, \ldots, v_{m})$.
\end{question}

\begin{solution}
	To prove that $\spn(w_1, \ldots, w_{m}) = \spn(v_1, \ldots, v_{m})$, we must show that $\spn(w_1, \ldots, w_{m}) \subseteq \spn(v_1, \ldots, v_{m})$ and $\spn(v_1, \ldots, v_{n})$.

	\begin{step}
		Prove $\spn(w_1, \ldots, w_{m}) \subseteq \spn(v_1, \ldots, v_{m})$.
	\end{step}

	Let $u \in \spn(w_1, \ldots, w_{m})$. By definition of span, we know that there exists $\lambda_1, \ldots, \lambda_m \in \mathbb{F}$ such that:

	\begin{gather*}
		\lambda_1{w_1} + \ldots + \lambda_m{w_m} = u \\
	\end{gather*}

	However, we know that $w_k = v_1 + \ldots + v_k$ by its definition. Therefore, $\lambda_k{w_k} = \lambda_k(v_1 + \ldots + v_{k}) = \lambda_k{v_1} + \ldots + \lambda_k{v_k}$. Therefore, we have that $\lambda_k{w_k} + \lambda_{k + 1}{w_{k + 1}} + \ldots + \lambda_m{w_m} = (\lambda_k + \lambda_{k + 1} + \ldots \lambda_{m})v_k + (\lambda_{k + 1} + \ldots + \lambda_{m}){v_{k + 1}} + \ldots \lambda_m{v_m}$ and hence:

	\begin{alignat*}{2}
		\lambda_1{v_1} + \ldots + \lambda_m(v_1 + \ldots + v_{m}) &= u && \implies \\
		\lambda_1{v_1} + \ldots + \lambda_m{v_1} + \ldots + \lambda_m{v_m} &= u && \implies \\
		(\lambda_1 + \ldots + \lambda_{m})v_1 + \ldots + \lambda_m{v_m} &= u \\
	\end{alignat*}

	Thus, since $(\lambda_k + \ldots + \lambda_{m}) \in \mathbb{F}$, letting $\alpha_k \in \mathbb{F}$ such that $\alpha_k = \sum_{i = k}^{m}\lambda_i$, we have that:

	\begin{gather*}
		\alpha_1{v_1} + \ldots + \alpha_m{v_m} = u \\
	\end{gather*}

	And thus $u \in \spn(v_1, \ldots, v_{m})$. Therefore, $\spn(w_1, \ldots w_{m}) \subseteq \spn(v_1, \ldots, v_{m})$.

	\begin{step}
		Prove $\spn(v_1, \ldots, v_k) \subseteq \spn(w_1, \ldots, w_k)$.
	\end{step}

	Let $u \in \spn(v_1, \ldots v_m)$. By definition of span, we know that there exists $\lambda_1, \ldots, \lambda_m \in \mathbb{F}$ such that:

	\begin{gather*}
		\lambda_1{v_1} + \ldots + \lambda_{m}{v_m} = u \\
	\end{gather*}

	Now, let $\alpha_k \in \mathbb{F}$ such that $\alpha_m = \lambda_m, \alpha_{m - 1} = \lambda_{m - 1} - \lambda_{m}, \ldots, \alpha_1 = \lambda_1 - \lambda_2 - \ldots \lambda_m$. By definition of $w_k$, we have that:

	\begin{align*}
		\alpha_1{w_1} + \ldots + \alpha_m{w_m} &= \alpha_1{v_1} + \ldots + \alpha_m(v_1 + \ldots + v_m) \\
		&= \alpha_1{v_1} + \ldots + \alpha_m{v_1} + \ldots + \alpha_m{v_m} \\
	\end{align*}

	And by our construction of $\alpha_k$, we have that:

	%TODO: Tighten up proof a bit

	\begin{align*}
		\alpha_1{v_1} + \ldots + \alpha_m{v_1} + \ldots + \alpha_m{v_m} &= (\lambda_1 - \ldots - \lambda_m)v_{1} + \ldots + \lambda_{m}{v_1} + \ldots + \lambda_{m}{v_m} \\
		&= \lambda_1{v_1} - \ldots - \lambda_m{v_1} + \ldots + \lambda_m{v_1} + \ldots + \lambda_m{v_m} \\
		&= \lambda_1{v_1} + \ldots + \lambda_m{v_m} \\
		&= u \\
	\end{align*}

	And thus $u \in \spn(w_1, \ldots, w_m)$. Therefore, $\spn(v_1, \ldots, v_m) \subseteq \spn(w_1, \ldots, w_2)$. \\

	Therefore, since $\spn(w_1, \ldots, w_m) \subseteq \spn(v_1, \ldots, v_m)$ and $\spn(v_1, \ldots, v_m) \subseteq \spn(w_1, \ldots, w_m)$, $\spn(v_1, \ldots, v_m) = \spn(w_1, \ldots, w_m)$.
\end{solution}

\begin{question}
	Prove or give a counterexample: $v_1, \ldots, v_m$ is a linearly independent list of vectors in $V$, then $5v_1 - 4v_2, v_2, \ldots, v_m$ is linearly independent.
\end{question}

\begin{solution}
	Let $\lambda_1, \ldots, \lambda_m \in \mathbb{F}$. We can see that:

	\begin{alignat*}{2}
		\lambda_1(5v_1 - 4v_2) + \lambda_2{v_2} + \ldots + \lambda_m{v_m} &= 0 &&\implies \\
		5\lambda_1{v_1} - 4\lambda_1{v_2} + \lambda_2{v_2} + \ldots + \lambda_m{v_m} &= 0 &&\implies \\
		5\lambda_1{v_1} + (\lambda_2 - 4\lambda_1)v_2 + \ldots + \lambda_m{v_m} &= 0 && \\
	\end{alignat*}

	Now since $v_1, \ldots, v_m$ are linearly independent, we know that this implies $5\lambda_1 = (\lambda_2 - 4\lambda_1) = \ldots = \lambda_m = 0$. Thus, if $5\lambda_1 = 0$, then $\lambda_1 = 0$, and if $\lambda_1 = 0$, then $\lambda_2 - 4\lambda_1 = 0$ implies $\lambda_2 = 0$. Thus, we have that:

	\begin{gather*}
		\lambda_1(5v_1 - 4v_2) + \lambda_2{v_2} + \ldots + \lambda_m{v_m} = 0
	\end{gather*}

	Implies $\lambda_1 = \ldots = \lambda_m = 0$ and hence $5v_1 - 4v_2, v_2 \ldots, v_m$ is linearly independent.
\end{solution}

\begin{question}
	Let $v_1, \ldots, v_m$ be a linearly independent list of vectors in $V$, and let $w \in V$. Show that if the list $v_1, \ldots, v_m, w$ is a linearly independent list if and only if $w \notin \spn(v_1, \ldots, v_m)$.
\end{question}

\begin{solution}
	We will prove this by showing the forward and reverse directions separately.

	\begin{step}
		Forward direction.
	\end{step}

	Assume, for the purpose of contradiction, that $v_1, \ldots, v_m, w$ is a linearly independent list but that $w \in \spn(v_1, \ldots, v_m)$. Then by definition of span, we know there exists $\lambda_1, \ldots, \lambda_m \in \mathbb{F}$ such that:

	\begin{gather*}
		\lambda_1{v_1} + \ldots + \lambda_m{v_m} = w \\
	\end{gather*}

	However, note that this would imply that:

	\begin{gather*}
		\lambda_1{v_1} + \ldots + \lambda_m{v_m} + (-w) = 0 \\
	\end{gather*}

	Which implies that $v_1, \ldots, v_m, w$ cannot be a linearly independent list of vectors since $w$ has a non-zero coefficient. Therefore, by contradiction, $v_1, \ldots, v_m, w$ being a linearly independent list implies that $w \notin \spn(v_1, \ldots, v_m)$.

	\begin{step}
		Reverse direction
	\end{step}

	Once again, assume for the purposes of contradiction that $w \notin \spn(v_1, \ldots, v_m)$ but that both $v_1, \ldots, v_m$ and $v_1, \ldots, v_m, w$ are linearly dependent lists of vectors. Then, we know there exists $\lambda_1, \ldots, \lambda_m, \lambda_{m + 1} \in \mathbb{F}$ where some $\lambda_k \neq 0$ such that:

	\begin{gather*}
		\lambda_1{v_1} + \ldots + \lambda_m{v_m} + \lambda_{m + 1}w = 0 \\
	\end{gather*}

	However, this would imply that:

	\begin{gather*}
		\lambda_1{v_1} + \ldots + \lambda_m{v_m} = -\lambda_{m + 1}w \\
	\end{gather*}

	Now, if $\lambda_{m + 1} \neq 0$, then we have that:

	\begin{gather*}
		-\frac{\lambda_1}{\lambda_{m + 1}}{v_1} + \ldots - \frac{\lambda_m}{\lambda_{m + 1}}v_m = w \\
	\end{gather*}

	And thus, $w \in \spn(v_1, \ldots, v_m)$ which would be a contradiction. If $\lambda_{m + 1} = 0$, then we would have that:

	\begin{alignat*}{2}
		\lambda_1{v_1} + \ldots + \lambda_m{v_m} &= -\lambda_{m + 1}w &&\implies \\
		\lambda_1{v_1} + \ldots + \lambda_m{v_m} &= 0 \\
	\end{alignat*}

	And since at least one $\lambda_k \neq 0$, that would imply that $v_1, \ldots, v_m$ is not a linearly independent list of vectors, which would also be a contradiction. Therefore, $v_1, \ldots, v_m$ being an independent list of vectors and $w \notin \spn(v_1, \ldots, v_m)$ implies that $v_1, \ldots, v_m, w$ is a linearly independent list of vectors. \\

	Thus, for a linearly independent list of vectors $v_1, \ldots, v_m$, we have that $v_1, \ldots, v_m, w$ is also a linearly independent list of vectors if and only if $w \notin \spn(v_1, \ldots, v_m)$.
\end{solution}

\begin{question}
	Let $\mathbb{F}$ be a field. For $n \in \mathbb{N}$, let:

	\begin{gather*}
		\mathbb{F}{[x]}_n \coloneq \set{p \in \mathbb{F}[x] \colon \degree(p) \le n}
	\end{gather*}
\end{question}

\begin{subquestion}
	Prove that $\mathbb{F}{[x]}_n$ is a subspace of $\mathbb{F}[x]$.
\end{subquestion}

\begin{solution}

\end{solution}

\begin{subquestion}
	Explain why there does not exist a list of six polynomials that is linearly independent in $\mathbb{F}{[x]}_4$
\end{subquestion}

\begin{solution}
	%TODO: Clean up a bit
	We know that any list of six polynomials cannot be a linearly independent list in $\mathbb{F}{[x]}_4$ since a list of five polynomials forms a basis of $\mathbb{F}{[x]}_4$, and hence any list of six vectors in $\mathbb{F}{[x]}_4$ cannot be linearly independent.
\end{solution}

\begin{subquestion}
	Explain why no list of four polynomials spans $\mathbb{F}{[x]}_4$.
\end{subquestion}

\begin{solution}
	%TODO: Clean up a bit
	We know that any list of four polynomials cannot span $\mathbb{F}{[x]}_4$ since a list of five polynomials forms a basis of $\mathbb{F}{[x]}_4$, and hence any list of four vectors cannot span $\mathbb{F}{[x]}_4$.
\end{solution}

\begin{question}
	Consider the set $F_3 \coloneq \set{0, 1, 2}$ together with the addition $a \oplus b \coloneq (a + b) \mod 3$ and the multiplication $a \otimes b \coloneq (ab) \mod 3$.
\end{question}

\begin{subquestion}
	Verify that $\oplus$ and $\otimes$ satisfy the field axioms; only submit here the proof of distributivity.
\end{subquestion}

\begin{solution}

\end{solution}

\begin{subquestion}
	List all functions in the vector space ${F_3}^{F_3}$ (how many are there?), and determine which of these are polynomial functions.
\end{subquestion}

\begin{solution}
	There are a total of 27 functions in the vector space ${F_3}^{F_3}$:

	\begin{alignat}{3}
		f(2) &= 0 \quad f(1) &= 0 \quad f(0) &= 0 \\
		f(2) &= 0 \quad f(1) &= 0 \quad f(0) &= 1 \\
		f(2) &= 0 \quad f(1) &= 0 \quad f(0) &= 2 \\
		f(2) &= 0 \quad f(1) &= 1 \quad f(0) &= 0 \\
		f(2) &= 0 \quad f(1) &= 1 \quad f(0) &= 1 \\
		f(2) &= 0 \quad f(1) &= 1 \quad f(0) &= 2 \\
		f(2) &= 0 \quad f(1) &= 2 \quad f(0) &= 0 \\
		f(2) &= 0 \quad f(1) &= 2 \quad f(0) &= 1 \\
		f(2) &= 0 \quad f(1) &= 2 \quad f(0) &= 2 \\
		f(2) &= 1 \quad f(1) &= 0 \quad f(0) &= 0 \\
		f(2) &= 1 \quad f(1) &= 0 \quad f(0) &= 1 \\
		f(2) &= 1 \quad f(1) &= 0 \quad f(0) &= 2 \\
		f(2) &= 1 \quad f(1) &= 1 \quad f(0) &= 0 \\
		f(2) &= 1 \quad f(1) &= 1 \quad f(0) &= 1 \\
		f(2) &= 1 \quad f(1) &= 1 \quad f(0) &= 2 \\
		f(2) &= 1 \quad f(1) &= 2 \quad f(0) &= 0 \\
		f(2) &= 1 \quad f(1) &= 2 \quad f(0) &= 1 \\
		f(2) &= 1 \quad f(1) &= 2 \quad f(0) &= 2 \\
		f(2) &= 2 \quad f(1) &= 0 \quad f(0) &= 0 \\
		f(2) &= 2 \quad f(1) &= 0 \quad f(0) &= 1 \\
		f(2) &= 2 \quad f(1) &= 0 \quad f(0) &= 2 \\
		f(2) &= 2 \quad f(1) &= 1 \quad f(0) &= 0 \\
		f(2) &= 2 \quad f(1) &= 1 \quad f(0) &= 1 \\
		f(2) &= 2 \quad f(1) &= 1 \quad f(0) &= 2 \\
		f(2) &= 2 \quad f(1) &= 2 \quad f(0) &= 0 \\
		f(2) &= 2 \quad f(1) &= 2 \quad f(0) &= 1 \\
		f(2) &= 2 \quad f(1) &= 2 \quad f(0) &= 2
	\end{alignat}

	$\;$ \\

	All of them, in fact, have a polynomial representation. Letting $f$ be any of the functions listed above where $f(0) = a, f(1) = b$, and $f(2) = c$, we can express $f(x)$ as:

	\begin{gather*}
		f(x) = \frac{a}{2}(x - 1)(x - 2) - bx(x - 2) + \frac{c}{2}x(x - 1) \\
	\end{gather*}

	Since:

	\begin{align*}
		f(0) &= \frac{a}{2}(0 - 1)(0 - 2) - b(0)(0 - 2) + \frac{c}{2}(0)(0 - 1) \\
		&= \frac{a}{2}(-1)(-2) \\
		&= \frac{a}{2}(2) \\
		&= a \\\\
		f(1) &= \frac{a}{2}(1 - 1)(1 - 2) - b(1)(1 - 2) + \frac{c}{2}(1)(1 - 1) \\
		&= -b(1)(-1) \\
		&= -b(-1) \\
		&= b \\\\
		f(2) &= \frac{a}{2}(2 - 1)(2 - 2) - b(2)(2 - 2) + \frac{c}{2}(2)(2 - 1) \\
		&= \frac{c}{2}(2)(1) \\
		&= \frac{c}{2}(2) \\
		&= c \\
	\end{align*}

	Thus, since such an $f$ is clearly polynomial and can be used to construct all the functions listed above, all the functions in the vector space ${F_3}^{F_3}$ are polynomial.
\end{solution}

\begin{subquestion}
	Give a linearly independent list of polynomial functions that spans $\mathcal{P}(F_3)$.
\end{subquestion}

\begin{solution}

\end{solution}









\end{document}
